\vfill\pagebreak
\section{Making a decision: \texttt{if} and \texttt{else}}
\label{sec:control:if}

\token{if} runs a block only when a condition holds. The condition is any
expression of type \token{bool}, the type of Section~\ref{sec:types:boolean},
and the block runs when it is \token{true}. When it is \token{false}, the block
is skipped, and the program continues after it.

\begin{lstlisting}[style=coloredverbatim, caption=A block run only when a condition holds, label=lst:control:if]
use std::io;

fn main() {
  let temperature = read!i32(ask-> "Temperature: ");
  if temperature > 30 {
    println("It is hot, drink some water.");
  }
  println("Have a nice day.");
}
\end{lstlisting}
\vspace{-10pt}%
\begin{lstlisting}[style=bashVerb, escapechar=@+]
@+\textcolor{teal!80}{alice@dev:\mtilde}@+$ ./main
Temperature: 35
It is hot, drink some water.
Have a nice day.
@+\textcolor{teal!80}{alice@dev:\mtilde}@+$ ./main
Temperature: 12
Have a nice day.
\end{lstlisting}

The last line is printed on both runs, because it is outside the block: the
\token{if} decides about line~6 only.

\subsection{Two paths: \texttt{else}}

\token{else}, written after the block of an \token{if}, adds a second block that
runs when the condition is \token{false}. One of the two blocks always runs, and
never both. The next listing tells even numbers from odd ones with the remainder
operator \tokennolst{\%}, and Figure~\ref{fig:control:if_else} draws the two paths it
can take.

\begin{lstlisting}[style=coloredverbatim, caption=Two branches, label=lst:control:if_else]
use std::io;

fn main() {
  let n = read!i32(ask-> "A number: ");
  if n % 2 == 0 {
    println(n, " is even.");
  } else {
    println(n, " is odd.");
  }
}
\end{lstlisting}

\begin{figure}[h]
  \centering
  \begin{adjustbox}{max width=\linewidth}
    \begin{tikzpicture}[node distance=6mm and 4mm]

      \node[cfterm] (start) {start of main};
      \node[cfstep, below=of start] (read) {let n = read!i32(...);};
      \node[cftest, below=of read] (test) {n \% 2 == 0};
      \node[cfstep] (then) at ($(test) + (-32mm, -16mm)$) {println(n, " is even.");};
      \node[cfstep] (else) at ($(test) + (32mm, -16mm)$) {println(n, " is odd.");};
      \node[cfterm] (end) at ($(test) + (0, -30mm)$) {end of main};

      \draw[cfflow] (start) -- (read);
      \draw[cfflow] (read) -- (test);
      \draw[cfflow] (test.west) -| node[cflab, above, pos=0.4] {true} (then.north);
      \draw[cfflow] (test.east) -| node[cflab, above, pos=0.4] {false} (else.north);
      \draw[cfflow] (then.south) |- (end.west);
      \draw[cfflow] (else.south) |- (end.east);
    \end{tikzpicture}
  \end{adjustbox}
  \caption{\label{fig:control:if_else} The two paths through
    Listing~\ref{lst:control:if_else}. Exactly one of the two branches runs,
    then both paths meet again.}
\end{figure}


\subsection{More than two paths: \texttt{else if}}

The block after \token{else} can be replaced by another \token{if}, which gives a
chain of conditions. They are tested from top to bottom, and the first one that
holds wins: its block runs, and the rest of the chain is skipped. The final
\token{else} catches every case that no condition matched.

\begin{lstlisting}[style=coloredverbatim, caption=A chain of conditions, label=lst:control:else_if]
use std::io;

fn main() {
  let score = read!i32(ask-> "Score (0 to 100): ");
  if score >= 90 {
    println("Grade A");
  } else if score >= 75 {
    println("Grade B");
  } else if score >= 50 {
    println("Grade C");
  } else {
    println("Grade F");
  }
}
\end{lstlisting}

The order does the work here. A score of 95 is also above 75 and above 50, but
the first test catches it, so it gets an A and nothing else. Written in the
opposite order, starting with \token{score >= 50}, the chain would give a C to
everyone who passes.

The final \token{else} is optional, as it is after a single \token{if}. Without
it, a value that meets none of the conditions runs none of the blocks, and the
program simply continues after the chain. The next listing warns about the two
extremes of temperature and says nothing about the others.

\begin{lstlisting}[style=coloredverbatim, caption=A chain with no final \token{else}, label=lst:control:else_if_no_else]
use std::io;

fn main() {
  let temperature = read!i32(ask-> "Temperature: ");
  if temperature > 30 {
    println("It is hot, drink some water.");
  } else if temperature < 0 {
    println("It is freezing, wear a coat.");
  }
  println("Have a nice day.");
}
\end{lstlisting}
\vspace{-10pt}%
\begin{lstlisting}[style=bashVerb, escapechar=@+]
@+\textcolor{teal!80}{alice@dev:\mtilde}@+$ ./main
Temperature: -5
It is freezing, wear a coat.
Have a nice day.
@+\textcolor{teal!80}{alice@dev:\mtilde}@+$ ./main
Temperature: 15
Have a nice day.
\end{lstlisting}

\subsection{Choosing a value}

An \token{if} with an \token{else} is also an expression, whose value is the
value of the block that ran. Section~\ref{sec:control:blocks} gave each block a
value, so a decision between two values can be written in one line.

\begin{lstlisting}[style=coloredverbatim, caption=An \token{if} used as a value, label=lst:control:if_value]
use std::io;

fn main() {
  let n = read!i32(ask-> "A number: ");
  let parity = if n % 2 == 0 { "even" } else { "odd" };
  println(n, " is ", parity, ".");
}
\end{lstlisting}

For this to make sense, both blocks must produce a value, of the same type.
Without an \token{else}, the program would have no value to give
\token{parity} when \token{n} is odd, and the compiler refuses it:
\errmsg{E4081}{the if condition has no else value, so it must produce a void
  value not a mut [c8]}.

\begin{lstlisting}[style=coloredverbatimError, escapechar=@, caption=An \token{if} with no \token{else} has no value]
use std::io;

fn main() {
  let n = read!i32(ask-> "A number: ");
  let parity = @\hb{if n \% 2 == 0 \{ "even" \}}@;
  println(n, " is ", parity, ".");
}
\end{lstlisting}

The same rule rejects \token{if n > 0 \{ 1 \} else \{ "no" \}}, whose two
blocks produce an \token{i32} and a string: the type of a variable is fixed when
it is declared (Section~\ref{sec:types:variable_declaration}), so it cannot hold
one or the other depending on the run.

\subsection{Conditions the compiler refuses}

A condition must be a \token{bool}. Some languages accept a number and treat
zero as false; Ymir does not, and \token{if n \{} is an error,
\errmsg{E4095}{incompatible types i32 and bool}. The test has to be spelled out: \token{if n != 0 \{}.

A condition must also be unknown at compile time. In the next listing,
\token{limit} is a constant equal to 30, so the compiler knows the test is
always \token{true}: the \token{if} decides nothing, and the compiler says so
rather than let it pass.

\begin{lstlisting}[style=coloredverbatimError, escapechar=@, caption=A decision the compiler can already make]
use std::io;

fn main() {
  let limit = 30;
  if @\hb{limit > 20}@ {
    println("The limit is above 20.");
  }
}
\end{lstlisting}

The message is \errmsg{E4268}{conditional test is always evaluated to true,
  branch is always entered}. A condition that is always false gets the matching
message, E4266, since its
block could never run. A condition therefore has to depend on a value that the
compiler cannot know. In this section, that is always a value typed by the user
or drawn at random, with \token{read} and \token{uniform}, from
Section~\ref{sec:control:runtime_values}. The loops of the next sections add a
third kind: a variable that the program itself changes as it runs.

\notebox{The block of an \token{if} or of an \token{else} that holds a single
  instruction can be written without its braces: \token{if n < 0
  println("negative");} is valid. Not every block allows it. Apart from these
  two, only the \token{for} of Section~\ref{sec:control:for} and the arms of a
  \token{match} do; the body of a function, of a \token{loop} or of a
  \token{while} always needs its braces. The book writes them everywhere,
  because adding a second instruction to an \token{if} without braces silently
  puts it outside the \token{if}.}
